\subsection{优点：两角约束与可退化场模型}

同片双角应共同解释条纹，同时保留各自的响应差异。

共享厚度、逐角响应使条纹与背景分开解释。图\ref{fig:q2-residual}与表\ref{tab:q2-results}保留连续留段对照，避免逐角求厚度后平均。

场模型也能相互核对：直接取模平方与交叉项展开的最大差为$1.7\times10^{-16}$（反射率比例），图\ref{fig:field-expansion}验证了两种强度表达的一致性。
% src:求解/问题1/结果/模型验证.json:场展开一致性最大绝对差_比例

完整往返场沿用同一首返回场，往返振幅比趋零即退回两束。图\ref{fig:roundtrip}展示两模型的代数联系；分别估参的厚度差仍包含光学参数补偿。

测厚与高阶判断分开，图\ref{fig:20}与表\ref{tab:08}保留材料处理的区别，测厚不以检出高阶为前提。厚度与折射率的补偿、外部标定缺失，则是下节讨论的局限。
